% Generated draft: requires mathematical and source review.
\documentclass[11pt]{article}
\usepackage[margin=25mm]{geometry}
\usepackage{fontspec}
\setmainfont{Latin Modern Roman}
\usepackage{amsmath,amssymb,mathtools,unicode-math}
\setmathfont{Latin Modern Math}
\usepackage{enumitem,needspace,xurl,hyperref}
\setlength{\parindent}{0pt}
\setlength{\parskip}{5pt}
\setlength{\emergencystretch}{3em}
\newcommand{\N}{\mathbb{N}}
\newcommand{\Z}{\mathbb{Z}}
\newcommand{\Q}{\mathbb{Q}}
\newcommand{\R}{\mathbb{R}}
\newcommand{\C}{\mathbb{C}}
\newcommand{\F}{\mathbb{F}}
\newcommand{\A}{\mathbb{A}}
\newcommand{\PP}{\mathbb{P}}
\newcommand{\E}{\mathbb{E}}
\newcommand{\eps}{\varepsilon}
\newcommand{\dd}{\,\mathrm{d}}
\begin{document}
\section*{UnsolvedMath ID: 20002695}
\textbf{Source problem code:} catalogue entry
\subsection{Definitions and mathematical statement}
The source asks for a noncommutative Bakry--\'Emery criterion and an analogue of \(\Gamma_2\) for Dirichlet forms. The aid selects the algebraic operator-order formulation, not an assertion that all quantum curvature notions coincide.

Let \(P_t=e^{-tA}\) be a trace-symmetric unital completely positive Markov semigroup on a finite tracial von Neumann algebra, with positive \(L^2\)-generator \(A\). On a common \(*\)-algebra core where every product and iterated generator term is defined, put
\[
 \Gamma(x,y)=\tfrac12\bigl(A(x^*)y+x^*A(y)-A(x^*y)\bigr),
\]
\[
 \Gamma_2(x,y)=\tfrac12\bigl(\Gamma(x,Ay)+\Gamma(Ax,y)-A\Gamma(x,y)\bigr).
\]
For \(K\in\R\), \(N\in(0,\infty]\), the algebraic condition \(\mathrm{BE}(K,N)\) means
\[
 \Gamma_2(x,x)\succeq K\Gamma(x,x)+N^{-1}(Ax)^*(Ax)
\]
for all core elements, with \(1/\infty=0\). Complete \(\mathrm{CBE}(K,N)\) requires the same constants after every amplification \(A_r=\operatorname{id}_{M_r}\otimes A\). The operator-domain and matrix-level choices are essential parts of a precise criterion.

The aid's finite-dimensional model takes a nonidentity trace-preserving conditional expectation \(E:\mathcal M\to\mathcal B\) and \(A=\operatorname{id}-E\). For \(z=x-E(x)\),
\[
 \Gamma(x,x)=\tfrac12(z^*z+E(z^*z)),\qquad
 \Gamma_2(x,x)=\tfrac14z^*z+\tfrac34E(z^*z).
\]
Let \(C_E\) be the least constant with \(z^*z\preceq C_E E(z^*z)\) for all \(E(z)=0\), and let \(C_E^{\mathrm{cb}}\) be the corresponding least constant over all matrix amplifications. Under the stated finite-dimensional nontrivial centered setup, the aid gives the sharp frontier
\[
 K_E(N)=\frac12+\frac1{C_E+1}-\frac{2C_E}{N(C_E+1)},
\]
and the complete frontier with \(C_E\) replaced by \(C_E^{\mathrm{cb}}\). A condition holds exactly for \(K\) at most its frontier. For \(E=\operatorname{id}\), all terms vanish and this nontrivial-centered formula is not the appropriate restriction.

For \(\mathcal M=M_n(\C)\), \(n\ge2\), and \(E(x)=\tau_n(x)I\), the aid's exact centered constants are
\[
 C_E=n,\qquad C_E^{\mathrm{cb}}=n^2-1.
\]
Thus
\[
 K_{\mathrm{op}}(N)=\frac12+\frac1{n+1}-\frac{2n}{N(n+1)},\qquad
 K_{\mathrm{cb}}(N)=\frac12+\frac1{n^2}-\frac{2(n^2-1)}{Nn^2}.
\]
In particular, complete curvature \(K=1/2\) requires and permits \(N\ge2(n^2-1)\).

The complete sharpness witness uses \(k=n^2-1\) nonidentity Weyl unitaries \(W_1,\ldots,W_k\), orthonormal for \(\tau_n\) and of trace zero, and
\[
 Z=\sum_{j=1}^kE_{1j}\otimes W_j\in M_k\otimes M_n.
\]
Then \(E_k(Z)=0\), \(E_k(Z^*Z)=I_k\otimes I_n\), while \(\|Z^*Z\|=k\), attaining the complete centered constant. These sharp finite-dimensional bounds do not by themselves classify all noncommutative Dirichlet forms or prove every possible functional-inequality consequence of curvature.
\subsection{Short English statement}
The aid fixes an algebraic noncommutative \(\Gamma_2\) and distinguishes operator-order curvature from its complete matrix-level form. For the depolarizing generator on \(M_n\), it gives sharp curvature-dimension frontiers with centered constants \(n\) and \(n^2-1\), respectively.
\subsection{Sources}
\begin{itemize}
\item Local registry, permanent UnsolvedMath ID \(20002695\), exact public code \texttt{catalogue entry}; canonical location \texttt{aim-probability-notes.json}, \texttt{notes[136]}.
\item AIM, \emph{Free Analysis}, Question 35, section 0.11, ``Focus Group on Dirichlet Forms, from Classical to Quantum (day 5)''; supplied reference \url{https://aimath.org/WWN/freeanalysis/freeanalysis.pdf}. The verification field reports prior visual inspection of physical/PDF page 7. No PDF was accessed here. No proposer, original remarks, or literature entry is supplied.
\item Supplied research attempt 1: \texttt{partial\_result}, run and registry status \texttt{partially\_solved}, rigor label \texttt{proved}, low novelty confidence; classification \texttt{PARTIAL-PROGRESS}, created and updated 14 August 2026. The standard algebraic iterate and sharp centered-index calculations are supplied without primary bibliographic references.
\item The research aid explicitly says it is machine generated, not peer reviewed, and requires independent verification before citation.
\end{itemize}
\end{document}
